\scanpage{269}
\mspage{1}

Les localiseurs \unsure{fondamentaux}

1. \uline{Axiomes} $\mathrm{Loc}(1)$ \unsure{à|:} $\mathrm{Loc}(3)$ (et \unsure{variantes}) : \uline{localiseurs fondamentaux},

J'ai envie de reprendre l'exposé des axiomes
utiles pour un ``localiseur \unsure{fondamental}

\[
W\subset\mathrm{Fl}(\mathrm{Cat}),
\]

en me guidant cette fois \unsure{plus} sur la \unsure{signification}
géométrique des axiomes, \unsure{que} sur les
aléas $\pm$ techniques \unsure{qu'on|que je} rencontre \unsure{(ou ne rencontre pas)} pour les \unsure{certifier|vérifier}. \unsure{Je suivrai|Je suis} \unsure{donc}
dans cet esprit la présentation \unsure{de|du} $\unsure{\mathrm{XV}|\mathrm{X}}$ § 9.

\uline{Loc} 1 (\uline{saturation}) \quad a) $\mathrm{Iso}\subset W$

b) Si deux parmi $\unsure{f,g,gf|f,g,fg}$ dans $W$, aussi le \unsure{$3^{\mathrm{e}}$|$3^o$}

c) Si $gf$ et $fg$ \unsure{sont} dans $W$, $f$ et $g$ \unsure{aussi}.

\margin{\unsure{1|!}}Les conditions a) \inserted{et} b) tout \unsure{seuls|seul}, je vais les appeler
(pour un $M\unsure{\in|\subset}\mathrm{Fl}\,M$ général) \uline{saturation faible} du
localiseur. Les conditions a) b) c), \uline{saturation modérée}
(``mild saturation''). \unsure{Autres variantes|Autre variante}

\uline{Loc} $(1\,\mathrm{bis})$. Les conditions a) b) c) de Loc 1, plus

d) $W$ stable par facteurs directs.

C'est bien sûr impliqué par :

\uline{Loc} $(1\,\mathrm{ter})$ $W$ est \uline{fortement saturé}, i.e. si $\unsure{f\in\mathrm{Fl}\,\mathrm{Cat}|f\in\mathrm{Cat}}$ est
telle que $\mathrm{hot}_{W}(f)$ \unsure{iso}, alors $f\in W$.

\deleted{\unsure{2|L}} \uline{Loc} $\unsure{(2)}$ (\uline{Objet final}). Si $X$ dans $\mathrm{Cat}$ a un \unsure{obj.|objet} final,
alors $X$ est $W$-asphérique i.e. $X\longrightarrow e$ est dans $W$.

Parfois, cet axiome est conséquence de la variante
affaiblie suivante

\uline{Loc} $\unsure{(2_0)|(2_o)}$ (\unsure{Axiome} du segment) $\Delta^1$ est $W$-asphérique
% archive sheet number: 268
